% Vietnamese translation for OI Public Library
% Source-File: IOI中国国家候选队论文/2008/Day1/1.曹钦翔《数据结构的提炼与压缩》/数据结构的提炼与压缩.ppt
\documentclass[11pt]{article}
\usepackage{oipl-vn}
\usepackage{tikz}
\usetikzlibrary{arrows.meta,positioning}

\newcommand{\Oh}{\mathcal{O}}
\renewcommand{\figurename}{Hình}

\title{Bản trình chiếu: Tinh luyện và nén cấu trúc dữ liệu}
\author{Cao Qinxiang}
\date{2008}

\begin{document}
\maketitle

\origtitle{Tinh luyện và nén cấu trúc dữ liệu: bản trình chiếu}
\sourcepath{IOI China candidate team papers / 2008 / Day 1 / Cao Qinxiang / slides.ppt}

\section{Mạch chính của slide}

Bản trình chiếu đi cùng bài viết \emph{Tinh luyện và nén cấu trúc dữ liệu}.
Luồng văn bản trong tệp PowerPoint cũ chỉ trích xuất được một phần, nhưng các
tiêu đề và ví dụ chính khớp với bài Word: \textit{Ural 1568 Train Car
Sorting}, \textit{CEOI 2007 Necklaces}, \textit{Hide and Seek}, thống kê trên
cây, truy vấn trên đường trong trie, và hai ví dụ dùng cây DFS/BFS. Bản dịch
này giữ lại mạch trình bày của slide, không lặp toàn bộ chứng minh dài của bài
viết.

Thông điệp trung tâm: một cấu trúc dữ liệu tốt không nhất thiết lưu nhiều
thông tin hơn. Trái lại, ta thường thắng bằng cách xác định thông tin nào là
vô dụng, thông tin nào bị lặp lại, và cấu trúc nào có thể ép về một mô hình
đơn giản hơn.

\section{Tinh luyện: bỏ thông tin rỗng}

Ví dụ mở đầu là \textit{Train Car Sorting}. Bài toán có thể mô tả một hoán vị
bằng ``ma trận mẹ''. Nếu lưu cả ma trận thì chi phí trở thành bậc hai, nhưng
trong ma trận chỉ có \(n\) ô khác rỗng. Các ô bằng không không ảnh hưởng tới
thao tác đọc hàng/cột và không cần nằm trong cấu trúc dữ liệu.

Vì vậy slide nhấn mạnh bước tinh luyện:
\begin{itemize}
  \item tìm mô hình thuận tiện để chứng minh thuật toán;
  \item nhìn lại mô hình đó và xóa các phần không tham gia vào tính toán;
  \item lưu các đối tượng thật sự hoạt động, không lưu phần nền trống.
\end{itemize}
Với \textit{Train Car Sorting}, cách lưu này giảm bộ nhớ xuống \(\Oh(n)\) và
tổng thời gian tỷ lệ với kích thước lời giải xuất ra.

\section{Thu gọn: chia sẻ phần chung}

\textit{CEOI 2007 Necklaces} là ví dụ chính của thu gọn thông tin lặp lại. Mỗi
thao tác tạo một chuỗi mới từ chuỗi cũ bằng cách thêm hoặc xóa ở đầu trái hay
đầu phải. Nếu sao chép toàn bộ chuỗi sau mỗi thao tác, bộ nhớ có thể tăng tới
\(\Oh(n^2)\). Nhưng hai phiên bản liên tiếp chỉ khác một phần tử, nên phần
chung cần được chia sẻ.

Bài viết dùng cấu trúc left-right tree: một chuỗi được tách thành phần trái
và phần phải. Phần trái được biểu diễn bằng một đường trong trie trái, phần
phải bằng một đường trong trie phải. Một phiên bản chuỗi chỉ cần bốn con trỏ
tới hai đường này.

\begin{figure}[htbp]
\centering
\begin{tikzpicture}[node distance=1.0cm, every node/.style={font=\small}]
  \node[circle,draw] (lr) {L};
  \node[circle,draw,below left=of lr] (la) {a};
  \node[circle,draw,below=of lr] (lb) {b};
  \node[circle,draw] (rr) at (4,0) {R};
  \node[circle,draw,below=of rr] (rc) {c};
  \node[circle,draw,below right=of rr] (rd) {d};
  \draw[-{Latex[length=2mm]}] (la) -- (lr);
  \draw[-{Latex[length=2mm]}] (lb) -- (lr);
  \draw[-{Latex[length=2mm]}] (rr) -- (rc);
  \draw[-{Latex[length=2mm]}] (rr) -- (rd);
  \draw[thick,blue] (la) -- (lr);
  \draw[thick,red] (rr) -- (rd);
  \node[align=center] at (2,-1.8) {một chuỗi = đường trong trie trái\\ghép với đường trong trie phải};
\end{tikzpicture}
\caption{Sơ đồ hóa left-right tree: nhiều phiên bản chuỗi chia sẻ các đường trie chung.}
\end{figure}

Vấn đề còn lại là đi nhanh trên một đường trie, chẳng hạn tìm con kế tiếp của
một tổ tiên trên đường tới lá. Slide đưa về bảng nhảy tổ tiên giống LCA:
lưu độ sâu và các ``siêu cha'', rồi nhảy theo lũy thừa hai.

\section{Nén cây thành dãy}

\textit{Hide and Seek} minh họa việc nén cấu trúc cây thành chuỗi ngoặc DFS.
Khoảng cách giữa hai đỉnh có thể đọc từ đoạn ngoặc giữa hai lần xuất hiện:
sau khi triệt tiêu các cặp ngoặc khớp, số dấu đóng còn lại là số bước đi lên,
số dấu mở còn lại là số bước đi xuống.

Nhờ đó mỗi đoạn được mô tả bởi một cặp \((a,b)\), và phép ghép hai đoạn trở
thành một phép hợp nhất cục bộ. Cây đoạn trên chuỗi ngoặc lưu thêm các đại
lượng tốt nhất ở tiền tố/hậu tố, từ đó trả lời đường kính của các đỉnh đang
được đánh dấu và hỗ trợ đảo màu một đỉnh.

\section{Hai thứ tự DFS}

Trong bài thống kê trên cây, với mỗi đỉnh \(v\) cần đếm số hậu duệ có nhãn nhỏ
hơn \(v\). Thay vì duyệt lại toàn bộ cây con của từng đỉnh, slide dùng hai
thứ tự:
\begin{itemize}
  \item DFS thông thường;
  \item DFS sau khi đảo thứ tự con tại mỗi nút.
\end{itemize}
Giao của hai vùng ``sau \(v\)'' trong hai dãy này chính là phần hậu duệ cần
đếm, trừ các hiệu chỉnh tổ tiên. Fenwick tree xử lý số lượng nhỏ hơn trên dãy
và trên đường gốc-đỉnh, đưa bài toán về \(\Oh(n\log n)\).

\section{DFS và BFS trong đồ thị}

Phần cuối của slide chuyển từ cây sang đồ thị. Với cây DFS của đồ thị vô
hướng, không có cạnh ngang; mọi cạnh ngoài cây nối một đỉnh với tổ tiên của
nó. Tính chất này làm các bài liên quan tới cầu, khớp, hướng cạnh và điều kiện
liên thông trở nên dễ kiểm soát hơn.

Ngược lại, cây BFS giữ thông tin khoảng cách theo lớp. Trong bài chọn cây
khung có đường kính nhỏ, thử từng gốc BFS và xét độ sâu các lớp cho ta ứng
viên tốt; khi đường kính tối ưu lẻ, tâm có thể là một cạnh nên cần kiểm tra
thêm điều kiện hai phía của cạnh tâm.

\section{Kết luận của slide}

Ba động tác ``tinh luyện, nén, thu gọn'' thực chất đều nhằm giảm quy mô thông
tin phải xử lý. Sau khi ép cấu trúc phức tạp về dãy, trie dùng chung, chuỗi
ngoặc, cây DFS hoặc cây BFS, ta có thể dùng những công cụ quen thuộc như cây
đoạn, Fenwick tree và bảng nhảy tổ tiên. Đó là cách biến cấu trúc dữ liệu khó
thành cấu trúc nhỏ hơn, nhanh hơn và dễ cài đặt hơn.

\end{document}
